% ---------------------------------------------------------------------
% Preambulo: classe, pacotes e metadados do trabalho
% (Este arquivo é incluído por main.tex)
% ---------------------------------------------------------------------
\documentclass[
    12pt,
    oneside,
    sumario=tradicional,
    idioma=brazil,
    bib=alf,
    refback,
    numb=chap
]{unbtex}

% --------------------------------------------------------
% Pacotes usuais (adicione outros conforme necessidade)
% --------------------------------------------------------
\usepackage[utf8]{inputenc}
\usepackage[T1]{fontenc}
\usepackage{lmodern}
\usepackage{microtype}
\usepackage{graphicx}
\usepackage{listings}
\usepackage{amsmath,amssymb}
\usepackage{hyperref}

\hypersetup{
    colorlinks=true,
    linkcolor=blue,
    citecolor=blue,
    urlcolor=magenta
}

% --------------------------------------------------------
% Informações do trabalho (preencha conforme necessário)
% --------------------------------------------------------
\titulo{Provisionamento Dinâmico de Máquinas Virtuais em Nuvem Privada Acadêmica com OpenStack}
\autori[]{Kauan de Torres}{Eiras}
\orientador[Orientador]{Prof. Dr.}{Sébastien Roland Marie Joseph}{Rondineau}
\tipotrabalho{Trabalho de Conclusão de Curso}
\tipocurso{Engenharia de Software}
\instituicao[Universidade de Brasília]{Faculdade de Tecnologia}{Departamento de Engenharia de Software}
\local{Brasília}
\ano{2025}
\preambulo{Trabalho de Conclusão de Curso apresentado ao curso de Engenharia de Software da Universidade de Brasília como requisito parcial para obtenção do grau de Engenheiro de Software.}

% Cria índice remissivo se desejar
%\makeindex
